\par\Needspace{4\baselineskip}
\begin{center}
\SourceLeaf{025b}\pdfbookmark[2]{7.}{article-007}7.
\end{center}
Thus, at perihelion the true anomaly, the eccentric anomaly and the mean anomaly are $= 0$.
As the true anomaly then increases, the eccentric and mean anomalies also increase.
The eccentric anomaly remains smaller than the true anomaly and the mean anomaly smaller
than the eccentric anomaly until aphelion, where all three simultaneously become $= 180^0$.
From there to perihelion the eccentric anomaly is always greater than the true anomaly
and the mean anomaly greater than the eccentric anomaly, until all three become $= 360^0$
at perihelion or, equivalently, all are again $= 0$. In general, if the true anomaly $v$ corresponds to the eccentric anomaly $E$ and the
mean anomaly $M$, then it is clear that the true anomaly $360^0-v$ corresponds to the
eccentric anomaly $360^0-E$ and the mean anomaly $360^0-M$. The difference $v-M$ between the true and mean anomalies is called the
\emph{equation of centre}. It is therefore positive from perihelion to aphelion and
negative from aphelion to perihelion, and vanishes at perihelion and aphelion themselves.
Since $v$ and $M$ traverse a full circle from $0$ to $360^0$ in the same time, the time
of one revolution, also called the \emph{orbital period}, is obtained in days by
dividing $360^0$ by the daily motion $\frac{k\BareRoot{(1+\mu)}}{a^{\frac{3}{2}}}$.
It follows that, for different bodies revolving around the Sun, the squares of the
orbital periods are proportional to the cubes of the mean distances insofar as
their masses or, more precisely, the differences between their masses can be neglected.
